\chapter{Equations Différentielles}


\setlength{\parindent}{0pt}
\renewcommand{\labelitemi}{\textbullet} % Utiliser des points noirs (•)

% ==================================================================================================================================
% Définitions et propriétés

\section{Définitions et propriétés}

Nous allons chercher dans ce cours à définir proprement la notion d'équation différentielles, 
de solution d'une équation différentielle et voir quelques propriétés de ces deux objets. 

\subsection{Schéma Récapitulatif}

% ---- Ici le schéma du cours avec exemples 

\subsection{Définitions}

\begin{definition}[Equation Différentielle Scalaire]
    Soit $f : \Omega \longrightarrow \K$ où $ \Omega \subseteq \R \times \K^N$ avec $N \in \N$. 
    L'équation différentielle associée à $f$ est :
        \[ (E) : y^{(n)} = f(t,y,\dots,y^{(n-1)}) \] 
    On appelle cela une \textbf{équation différentielle d'ordre n scalaire}. 
    On peut appeler $f$ la fonction décrivant l'équation différentielle. 
\end{definition}

Définissons maintenant clairement le concept de "solution" d'une équation différentielle. 

\begin{definition}[Solution]
    Une fonction $y$ est appelée solution d'une équation différentielle $(E)$ d'ordre $n \in \N$ 
    définie par la fonction $f$ si :
    \begin{itemize}
        \item $y : I \longrightarrow \K$ est $n$ fois dérivable 
        \item $ \forall t \in I, \quad (t,y(t), \dots, y^{(n-1)}(t)) \in \Omega $ 
        \item $ \forall t \in I$, on a : 
            \[ \boxed{ y^{(n)}(t) = f(t,y(t), \dots, y^{(n-1)}(t)) } \] 
    \end{itemize}
\end{definition}

\begin{remark}
    Si on change l'intervalle de définition $I$ d'une solution d'une équation différentielle, 
    alors on change aussi la solution. Autrement dit, une solution est très dépendante de son intervalle 
    de définition. On peut donc parfois noter $(I,y)$ une solution $y$ définie sur un intervalle $I$. 
\end{remark}

\begin{definition}[Equation Différentielle Vectorielle]
    Soit $f : \Omega \longrightarrow \K^N$ avec $n \in \N^*$ et $\Omega \subseteq \R \times (\K^N)^N$. 
    L'équation différentielle ordinaire associée à $f$ est :
        \[ (E) : y^{(n)} = f(t,y,\dots,y^{(n-1)}) \] 
    On peut écrire cette équation différentielle sous forme de système d'équations différentielles 
    en décomposant les différentes composantes de l'équation différentielle tel que :
        \[
            \begin{cases}
                y_1^{(n)} = f_1(t,y_1, \dots, y_N^{(n-1)}) \\ 
                \vdots \quad  \quad \quad \quad \quad \quad \quad \vdots \\ 
                y_N^{(n)} = f_1(t,y_N, \dots, y_N^{(n-1)}) \\ 
            \end{cases} 
        \]
\end{definition}

Lors de l'étude d'un equation différentielle d'ordre $n$, un théorème que nous verrons plus tard nous permettra
de nous ramener à l'étude d'une équation différentielle d'ordre $1$. 
Il est donc important de bien comprendre ce qu'est une équation différentielle d'ordre $1$. 

\begin{remark}[Equation Differentielle d'ordre 1]
    Soit $f : \Omega \longrightarrow \K^N $ où $ \Omega \subseteq \R \times \K^N$. 
    L'équation différentielle associée à $f$ est :
        \[ (E) : y' = f(t,y) \] 
\end{remark}

\begin{example}
    Voyons quelques exemples d'équations différentielles : 
    \begin{itemize}
        \item $\K = \R, N = 1, n = 1, \Omega = \R \times \R$ et $f(t,x) = x $ 
            \[ (E) : y' = y \] 
        \item $\K = \R, N = 2, n = 1, \Omega = \R \times \R^2$
            \[ f : 
                \begin{cases}
                    \Omega \longrightarrow \R^2 \\ 
                    (t,x) \longmapsto Ax + B 
                \end{cases}
                \text{ où } B \in \R^2 \text{ et } A \in \mathcal{M}_2(\R) \] 
            \[ (E) : y' = Ay + B \] 
            \[ (E) : 
                \begin{cases}
                    y_1' = a_{1,1}y_1 + a_{1,2}y_2 + b_1 \\ 
                    y_2' = a_{2,1}y_1 + a_{2,2}y_2 = b_2
                \end{cases}
            \]
        \item $\K = \C, N = 1, n = 1, \Omega = \R \times \C$ 
            \[ f : 
                \begin{cases}
                    \Omega \longrightarrow \C \\ 
                    (t,z) \longmapsto z^2 + t 
                \end{cases}
                \quad \text{ et } \quad 
                (E) : y' = y^2 + t 
            \]
    \end{itemize}
\end{example}

\subsection{Premières Propriétés}

Lors de la résolution d'une équation différentielle, on remarque que on trouve la plupart du temps une infinité de 
solutions. Parfois, une unique solution serait plus pratique (ex : en physique). On va donc chercher à donner des 
conditions supplémentaires sur la fonction solution en imposant une de ses valeurs. 

\begin{definition}[Problème de Cauchy]
    Soit $(E)$ une équation différentielle ordinaire d'ordre 1, de fonction descriptive $f$ et de la forme $y' = f(t,x)$. 
    Résoudre un problème de Cauchy associé à $(E)$ consiste à trouver une solution de $(E)$ qui vérifie un condition initiale : 
    \[ (C) : y(t_0) = y_0 \quad \text{ où } (t_0,y_0) \in \Omega \] 
\end{definition}

\begin{proposition}[Exercice]
    Soit $f : \Omega \subset \R \times \K^N \longrightarrow \K^N$ et $(E)$ l'équation associée d'ordre 1 :
        \[ y' = f(t,y) \] 
    Alors : $y : I \longrightarrow \K^N$ est solution du problème de Cauchy :
        \[ (P) : 
            \begin{cases}
                (E) : y' = f(t,y) \\ 
                (C) : y(t_0) = y_0 
            \end{cases}
        \] 
    si et seulement si 
        \begin{itemize}
            \item $y \in \mathcal{C}^0(I)$ 
            \item $ \forall t \in I, \quad (t,y(t)) \in \Omega $ 
            \item $ \forall t \in I$, 
                \[ y(t) = y_0 + \int_{t_0}^{t} f(s,y(s)) \; ds \] 
        \end{itemize}
\end{proposition}

\begin{definition}[Solution Maximale]
    $y$ est une solution maximale de $(E)$ si il n'existe pas de prolongement strict de $y$ en une solution de $(E)$. 
\end{definition}

\begin{definition}[Solution Globale]
    Soit $(E) : y' = f(t,y)$ où $f : I \times \Omega \longrightarrow \K^N$. 
    On dit que $y$ est une solution globale de $(E)$ si $y$ est une solution de $(E)$ et qu'elle est définie sur $I$.
\end{definition}

\begin{remark}
    Une solution globale est maximale. La réciproque est fausse, voyons un exemple. 
        \[ (E) : y' = y^2 \quad \text{ et } \quad y_s :
            \begin{cases}
                A \longrightarrow \R \\ 
                t \longmapsto - \frac{1}{t}
            \end{cases}
            \text{où } f: 
            \begin{cases}
                \R \times \R \longrightarrow \R \\ 
                (t,x) \longmapsto x^2
            \end{cases}
            \quad \text{ et } \quad A = \R_+^* \] 
    $y_s$ n'est pas une solution globale car on ne peut pas la prolonger en une fonction continue solution de $(E)$. 
    En revanche, c'est une solution maximale. 
\end{remark}

\begin{theorem}[Existence et Unicité de solution maximale]
    Il existe une solution maximale d'un problème de Cauchy. 
\end{theorem}

% ==================================================================================================================================
% Première méthode de résolution

\section{Première méthode résolution}

A partir des outils que nous avons à notre porté, nous pouvons dès à présent déterminer la forme de l'ensemble des solution 
d'une équation différentielle scalaire homogène d'ordre 1. Soient $I$ un intervalle, $a : I \longrightarrow \K$ une fonction continue 
et $(E)$ une équation différentielle de la forme : 
    \[ (E) : y' = a(t)y = f(t,y) \quad \text{où} \quad f : 
        \begin{cases}
            I \times \K \longrightarrow \K \\ 
            (t,x) \longmapsto a(t)x
        \end{cases}
    \]
Soit le problème de Cauchy suivant : 
    \[ (P) : 
        \begin{cases}
            (E) : y' = a(t)y = f(t,y) \\ 
            y(t_0) = y_0
        \end{cases}
        \quad \text{où} \quad (t_0,y_0) \in I \times \K \] 
On cherche la forme générale de l'ensemble des solutions du problème de Cauchy. 

\subsection{Existence d'une solution}

Soit $y : I \longrightarrow \K, t \longmapsto e^{ \int_{t_0}^{t} a}$. Montrons que $y$ est bien une solution maximale du 
problème de Cauchy. 
\begin{itemize}
    \item Soit $t \in I$ alors $y(t) \in \K$ 
    \item $ \forall t \in I$, on a : $y'(t) = y_0 a(t) e^{ \int_{t_0}^{t} a}$ 
    Donc $y$ est bien une solution de $(E)$. De plus, $y$ est une solution globale donc elle est maximale. 
    \item Condition initiale : 
        \[ y(t_0) = y_0 e^{ \int_{t_0}^{t_0} a} = y_0 e^0 = y_0 \] 
    On a donc que $y$ est bien une solution maximale du problème de Cauchy énoncé. 
\end{itemize}

\subsection{Unicité de la solution}

Soit $u : J \subset I \longrightarrow \K$ une solution maximale de $(E)$ telle que $t_0 \in J$ et $u(t_0) = y_0$. 

Montrons que la fonction $v : t \longmapsto u(t) e^{ - \int_{t_0}^{t} a}$ est constante. 
$v$ est dérivable par composition/produit. On a donc $ \forall t \in J$, 
    \begin{align*}
        y'(t) &= u'(t) \times e^{ \int_{t_0}^{t} a} - u(t)a(t)e^{ \int_{t_0}^{t} a} \\ 
              &= e^{ \int_{t_0}^{t} a} (u'(t) - u(t)a(t))
    \end{align*}
or $u$ est un solution maximale de $(E)$ donc $ \forall t \in J, u'(t) = a(t)u(t) \iff u'(t) - u(t)a(t) = 0$ d'où :
    \[ y'(t) = 0 \] 
Donc $v$ est une fonction constante sur $J$. Déterminons cette constante. 
    \[ v(t_0) = u(t_0) e^{ - \int_{t_0}^{t_0} a} = u(t_0)e^0 = u(t_0) = y_0 \] 
Donc $ \forall t \in J, \; v(t) = y_0 $. 

Montrons que $u = y_{ |J}$ et que $I = J$. Soit $t \in J$, on a :
    \begin{align*}
        & v(t) = u(t) e^{ - \int_{t_0}^{t_0} a} \\ 
        \iff & u(t) = v(t) e^{ \int_{t_0}^{t_0} a} \\ 
        & u(t) = y_0 e^{ \int_{t_0}^{t_0} a} 
    \end{align*}
Donc $u = y_{ |J}$. De plus, $y$ est un prolongement de $u$ or $u$ est supposée maximale (i.e elle n'admet pas de prolongement solution)
donc $I = J$. 

\subsection{Conclusion}

Enfin, ces calculs nous permettent d'énoncer le résultat suivant : 

\begin{theorem}[Solutions d'une ED scalaire homogène d'ordre 1]
    Soient $I$ un intervalle, $a : I \longrightarrow \K$ une fonction continue et $f : I \times \K \longrightarrow \K$. 

    Soit $(E) : y' = a(t)y = f(t,y)$ l'équation différentielle homogène d'ordre 1 associée à $f$. 
    Soient $(t_0, y_0) \in I \times \K$, et $(P) : \{(E), y(t_0) = y_0\}$ un problème de Cauchy. 
    Alors les solutions maximales de $(P)$ sont exactement les fonctions de la forme : 
        \[ y : 
            \begin{cases}
                I \longrightarrow \K \\ 
                t \longmapsto y_0 \exp \left( \int_{t_0}^{t} a(s) \; ds \right)
            \end{cases}
        \] 
\end{theorem}




